\chapter{Results}
\label{ch:results}

This chapter presents the experimental findings as a connected investigation. It proceeds in the order in which the study actually developed, because that order is also the logical one: each experiment is motivated by a question that the previous result left open, and each ends by naming the question that motivates the next. The reader who follows the chapter in order will therefore see not only what was found but why each step was taken, which for a diagnostic result is the substance of the contribution. The investigation begins with a practical obstacle, the calibration of the samplers, which unexpectedly reveals the influence of embedding geometry. It proceeds to the annealing dynamics of the discrete sampler, where an apparent convergence turns out to be an artifact. It then confronts the Metropolis--Hastings correction, and finds that the same correction which rescues the discrete sampler destroys the continuous one, for a reason it measures precisely. This leads to the central result, the gradient fallacy, that following the true gradient is indistinguishable from following a random direction, on all five energy functions. Three diagnostic sections then supply the mechanism, through the linearization radius, the likelihood trap, and the anisotropy of the embedding space. The chapter closes by asking whether amortization escapes the failure, and finding that it does not, and by testing the plug-and-play constraint directly, and finding that it does not steer. Throughout, each reported number is annotated in the source of this document with a comment naming the results file from which it was drawn, so that every quantitative claim is traceable to the experiment that produced it, without cluttering the text the reader sees.

\section{A First Obstacle: Step-Size Calibration and Embedding Geometry}
\label{sec:results-stepsize}

The investigation began with a simple intention: to run the two samplers on GPT-2 Large using a step-size schedule that had already been made to work on Llama-3, and to compare their behaviour. The expectation was that the schedule would transfer, perhaps with minor adjustment, since both are Transformer language models and the samplers are the same. That expectation was wrong, and the way in which it was wrong is the first substantive finding.

On GPT-2 Large, with the Llama-tuned schedule, the samplers did not merely perform poorly; they did not move at all. The entropy of the discrete proposal collapsed to zero within a few steps and no token was ever changed. The sampler was not broken but trapped, and diagnosing why led directly to the geometry of the embedding space. The anisotropy analysis, reported in full in Section~\ref{sec:results-aniso}, measures the mean distance between neighbouring token embeddings as approximately $1.82$ in GPT-2 Large against approximately $0.59$ in Llama-3, and the mean pairwise distance as $2.77$ against $0.84$, a difference of roughly a factor of three in each case.
% SOURCE: results_diag/diag_anisotropy_gpt2sft.json (mean_nearest_neighbour_l2=1.8222,
% mean_pairwise_l2=2.7725) and results_diag/diag_anisotropy_llama3-8b.json (0.585, 0.835).
The reason for the trapped sampler follows immediately. A step size calibrated for the tighter geometry of Llama-3 is simply too small to carry the continuous state across a Voronoi cell boundary in the more dispersed geometry of GPT-2, so the projected token never changes and the sampler sits frozen in its initial cell. Recovering the correct scale required an oracle sweep over candidate step sizes, which yielded a schedule running from $\epsilon_{\text{start}} = 10.5$ to $\epsilon_{\text{end}} = 0.1$ for the GPT-2 family, and this schedule is used throughout the rest of the study.
% SOURCE: schedule established via the *.dls.*.oracle.s50 runs in results_gpt2_v2/,
% e.g. gpt2-large.dls.policy.nomh.gn.oracle.s50.json.

This obstacle is worth reporting rather than quietly fixing, for two reasons. The immediate practical lesson is that there is no universal step-size schedule for these methods; the correct scale is a property of the specific model's embedding geometry and must be rediscovered for each model. That alone qualifies the plug-and-play framing, since a method that must be recalibrated to each model's geometry before it will move at all is not, in the strong sense, model-agnostic. But the obstacle also raises a sharper question, and it is this question that motivates the ablation at the heart of the thesis. Once the step size is calibrated so that the sampler does move, does it move \emph{toward better sequences}, guided by the gradient, or does it merely move? Calibration ensures motion; it does not ensure that the motion is informed. The oracle configuration, which selects the ideal step size at each iteration using knowledge of the ground-truth token, was introduced precisely to separate these possibilities, because it establishes the best any schedule could do and thereby distinguishes a failure of calibration from a failure of guidance. This question is taken up in earnest in Section~\ref{sec:results-fallacy}, and answering it consumes the rest of the chapter. This section, meanwhile, has answered the first half of RQ1: theoretically faithful Langevin dynamics cannot be applied off the shelf, because the step size that makes the sampler move at all is model-specific.

\section{An Apparent Convergence: Annealing Dynamics and the Quenching Effect}
\label{sec:results-quench}

With the step size calibrated, the discrete sampler moves, and the next question is how it behaves over the course of a run. Run without the Metropolis--Hastings correction, it exhibits a striking and suggestive pattern: its metrics remain noisy and unimproving for the majority of the annealing schedule and then, near the end, drop abruptly to much better values. The natural reading of this pattern, and the one we initially entertained, is that the sampler has finally converged: that after a long period of exploration it has located a good sequence and settled into it. If that reading were correct, it would be modest good news for the method.

A controlled test refutes the reading, and the design of the test is the point. If the late drop marked genuine convergence, it would occur when the model had found a good sequence, and its timing would be governed by the difficulty of the problem, not by the schedule. If instead the drop were an artifact of the annealing schedule, its timing would track the schedule: lengthen the schedule, and the drop would move. We therefore extended the schedule from $50$ steps to $100$ and observed where the drop occurred. It moved, from around step $40$ to around step $85$, tracking the schedule rather than remaining fixed at the point where a good sequence might have been found.
% SOURCE: results_gpt2_v2/gpt2-large.dls.policy.nomh.gn.free.s50.csv (drop near step 40)
% vs gpt2-large.dls.policy.nomh.gn.free.s100.csv (drop near step 85); per-step
% trajectories in figures_gpt2/*dls*s50*trajectories and *s100*trajectories.
The drop is therefore not convergence but what we term a \textbf{quenching effect}, and Figures~\ref{fig:dls-traj-50} and~\ref{fig:dls-traj-100} show it at the two schedule lengths.

The mechanism is a direct consequence of the annealing schedule interacting with the absence of the correction. Early in the schedule the step size is large and the injected Gaussian noise dominates the update, so the sampler, lacking any mechanism to reject a bad proposal, bounces stochastically across the vocabulary, exhibiting high entropy and poor contextual fit. As the schedule anneals the step size toward zero, the noise vanishes, and the update degenerates into deterministic gradient descent, which collapses greedily into whatever local optimum happens to be nearest. The abrupt late-schedule drop is the signature of this collapse. The name is borrowed from metallurgy, where quenching denotes a rapid cooling that freezes a material into whatever configuration it happened to occupy at the moment the cooling began, and the analogy is exact: removing the noise freezes the sampler into its current local optimum rather than allowing it to continue exploring toward a better one. The drop looks like convergence because the metrics improve, but it is the improvement of a system that has stopped moving, not one that has found the right place.

\begin{figure}[t]
    \centering
    \includegraphics[width=0.9\textwidth]{figures/gpt2-large.dls.gn.free.s50_new_trajectories.png}
    % SOURCE: figures_gpt2/gpt2-large.dls.gn.free.s50_new_trajectories.png
    \caption{Discrete Langevin Sampler trajectories on GPT-2 Large over a 50-step schedule, averaged across sequences. The metrics remain unimproving for most of the run and drop abruptly near the end as the noise schedule approaches zero, the signature of quenching rather than convergence.}
    \label{fig:dls-traj-50}
\end{figure}

\begin{figure}[t]
    \centering
    \includegraphics[width=0.9\textwidth]{figures/gpt2-large.dls.gn.free.s100_new_trajectories.png}
    % SOURCE: figures_gpt2/gpt2-large.dls.gn.free.s100_new_trajectories.png
    \caption{The same sampler over a 100-step schedule. The abrupt drop has moved from around step 40 to around step 85, tracking the annealing schedule rather than remaining fixed. This confirms that the drop is an artifact of cooling and not evidence of model convergence.}
    \label{fig:dls-traj-100}
\end{figure}

There is a revealing contrast when the Metropolis--Hastings correction is enabled: the quenching pattern does not occur. The chain converges smoothly and early, within the first twenty steps, and remains stable thereafter, maintaining a low divergence for the rest of the schedule.
% SOURCE: compare gpt2-large.dls.policy.mh.gn.free.s50 (final KL 6.541, stable)
% against the nomh counterpart (final KL 9.499) in final_kl_by_model.csv.
In the discrete sampler, then, the theoretically motivated correction is also the empirically better choice, and this is worth marking clearly precisely because it is the exception. It is the one place in the entire study where making the sampler correct also makes it work. That observation naturally raises the next question, and it is a question the correction itself poses: if the correction rescues the discrete sampler, what does it do to the continuous one, whose landscape, as Section~\ref{sec:bg-samplers} warned, violates the smoothness the correction assumes? The next section takes up that question, and the answer is the opposite of what the discrete case might lead one to expect.

\section{A Correction That Breaks the Sampler: The Metropolis--Hastings Breakdown in Continuous Space}
\label{sec:results-mh}

The result of the previous section, that the Metropolis--Hastings correction improves the discrete sampler, sets up a natural expectation: that the correction, being theoretically motivated, ought to improve the continuous sampler too, or at least not harm it. Section~\ref{sec:bg-mh} gave a reason to doubt this expectation, namely that the continuous sampler's energy is defined through a projection whose drift is discontinuous, violating the Lipschitz condition the correction relies on. The experiment resolves the tension decisively and in favour of the doubt, and it produces what is, in our assessment, the single most striking result of the sampler study.

Applied to the continuous sampler, the correction is not merely unhelpful; it is catastrophic. The sampler flatlines, rejecting essentially every proposal that would change anything, and it does so in a way the trace experiment quantifies precisely. Over $10{,}000$ logged proposals from the continuous sampler with gradient normalization disabled and the correction enabled, the overall acceptance rate is $2.6\%$. That figure alone would be unremarkable, but its decomposition by whether the proposal crossed a Voronoi cell boundary is not. The acceptance rate for proposals that \emph{stayed} within a cell, and therefore changed no token, is $0.03\%$, while the acceptance rate for proposals that \emph{crossed} a boundary, and therefore actually did something, is $3.7\%$.
% SOURCE: results_diag/traces_gpt2sft.json, config cls_policy_gnoff_mh:
% accept_rate_overall=0.0261, accept_rate_stayed=0.000344, accept_rate_crossed=0.03665.
Both rates are tiny, but the structure they reveal is the point: almost every accepted move is one that leaves the projected token unchanged. The set of moves the correction accepts and the set of moves that accomplish anything are very nearly disjoint. For comparison, the discrete sampler under the correction accepts within-cell moves with probability $1.0$, because such moves are trivially reversible, and accepts boundary-crossing moves at a healthy rate, which is exactly why it works where the continuous sampler does not.
% SOURCE: results_diag/traces_gpt2sft.json, config dls_policy_gn_mh:
% accept_rate_stayed=1.0, accept_rate_crossed=0.0927, accept_rate_overall=0.1213.

The question this raises is the one the acceptance experiment was built to answer: is the correction rejecting these boundary-crossing moves because the moves are bad, or because the correction cannot verify their reversibility? The two possibilities correspond to the two terms of the acceptance ratio, and the decomposition separates them cleanly. For proposals that crossed a boundary, the mean \emph{target} log-ratio is $+4.60$. That is positive, which means these moves, on average, \emph{improve} the sequence probability: the distribution being sampled would welcome them. The moves are good. Yet the mean \emph{proposal} log-ratio for those same moves is $-1325$.
% SOURCE: results_diag/traces_gpt2sft.json, config cls_policy_gnoff_mh:
% mean_log_target_ratio_crossed=4.6035, mean_log_proposal_ratio_crossed=-1325.44.
The reverse proposal assigns effectively zero density to the return move, and this term alone, dwarfing the favourable target term by three orders of magnitude, drives the acceptance probability to zero. Figure~\ref{fig:mh-decomp} shows the two distributions side by side: the target term is centred near zero and frequently positive for boundary-crossing moves, while the proposal term for exactly those moves is enormous and negative. The rejection is driven by the proposal term, not the target term.

\begin{figure}[t]
    \centering
    \includegraphics[width=\textwidth]{figures/fig_mh_decomposition.pdf}
    % SOURCE: figures/fig_mh_decomposition.pdf produced by analyze_mh.py from traces_gpt2sft_mh.csv
    \caption{Decomposition of the Metropolis--Hastings acceptance ratio in the continuous sampler, for proposals that stay within a Voronoi cell against those that cross a boundary. The target term (left) is well-behaved and often favourable for boundary-crossing moves; the proposal term (right) collapses to large negative values for exactly those moves. The rejection is driven by the proposal term, not the target term.}
    \label{fig:mh-decomp}
\end{figure}

This is the measured form of the Lipschitz argument of Section~\ref{sec:bg-mh}, and it confirms the theory of \citet{roberts1996exponential} in the breach: the Metropolis-adjusted Langevin algorithm requires a Lipschitz-continuous drift, and the continuous sampler's energy does not have one. When a proposal crosses a cell boundary, the reverse proposal mean is computed from the drift on the far side of the boundary, which, because the drift is discontinuous there, points somewhere unrelated to the direction back toward the original state. The original state therefore falls deep in the tail of the reverse Gaussian, the reverse density vanishes, and the correction, doing exactly what it is designed to do, rejects the move as irreversible. The correction is not malfunctioning; the landscape is violating the assumption on which the correction depends. The consequence is both precise and damaging: the continuous sampler only appears to function when the correction is omitted, that is, when it is run as a biased optimizer that ignores reversibility. Making it a correct sampler paralyzes it entirely. The same behaviour is observed on both GPT-2 Large and Llama-3, so it is a property of the problem class rather than of a single model, and it is best summarized as a fundamental incompatibility between continuous Markov chain Monte Carlo and a discrete text representation accessed by projection. This answers RQ2 for the continuous sampler: making it theoretically correct does not make it work, it does the opposite, and the reason is now measured rather than conjectured.

The two sampler results so far, the quenching of the discrete sampler and the breakdown of the continuous one, concern how the samplers \emph{move}. Neither has yet asked the more fundamental question that underlies both: whether the gradient they follow, when they do move, carries any useful information at all. That question is the central one, and the next section confronts it directly.

\section{A Spatial View of the Breakdown: Trajectory Analysis}
\label{sec:results-traj}

The acceptance statistics of the previous section describe the breakdown in terms of proposals accepted and rejected. A complementary and more concrete view comes from watching where the sampler actually goes, which is what the trajectory experiment records: for a handful of sequences it logs the full continuous state at every step, along with the distance from that state to the nearest token embedding and the identity of the Voronoi cell the state occupies. Reducing these trajectories to two simple summaries makes the breakdown visible as a matter of geometry rather than of acceptance ratios.

The first summary is the number of distinct Voronoi cells a run visits over its fifty steps, which is a direct count of how many times the projected token actually changed. The continuous sampler with the Metropolis--Hastings correction enabled visits, on average, exactly one cell across the entire run.
% SOURCE: traces_gpt2sft_traj.npz, cls_policy_gnoff_mh__cell_id: mean unique cells = 1.0.
It never leaves the cell it started in. This is the same paralysis reported in Section~\ref{sec:results-mh}, now seen from the side: a sampler that accepts only moves which change nothing is, spatially, a sampler that never leaves its initial cell, and the two descriptions are the same fact. By contrast, the continuous sampler with the correction disabled visits an average of $46.2$ of a possible $50$ cells, changing its token at almost every step.
% SOURCE: traces_gpt2sft_traj.npz, cls_policy_gnoff_nomh__cell_id: mean unique cells = 46.2.
It moves freely precisely because, without the correction, nothing checks whether its moves are reversible.

The second summary is the distance from the continuous state to the nearest token embedding, which measures how far the sampler has drifted from the manifold of actual tokens. Here the contrast between the two samplers is at its sharpest and it exposes a difficulty that the discrete sampler avoids by construction. The correction-free continuous sampler, in the course of forcing its way across cell boundaries, drifts as far as $588$ units from the nearest token embedding, ending its run an average of $17$ units away.
% SOURCE: traces_gpt2sft_traj.npz, cls_policy_gnoff_nomh__dist_to_manifold: max=588.2, end mean=17.3.
That is far outside the region where any token lives, given that the mean distance between neighbouring tokens is under two units, so the continuous state spends much of its run in the empty space between tokens, where the energy is defined only by an increasingly distant projection. The discrete sampler, by contrast, sits exactly on the token manifold at every step, at distance zero, because its state is always a genuine token and it never enters the continuous interstice at all.
% SOURCE: traces_gpt2sft_traj.npz, dls_policy_gn_mh__dist_to_manifold: identically 0.0.
It visits an average of seven cells over the run, moving deliberately from token to token rather than either freezing or drifting.
% SOURCE: traces_gpt2sft_traj.npz, dls_policy_gn_mh__cell_id: mean unique cells = 7.0.

These three trajectories, the frozen continuous sampler under the correction, the freely drifting continuous sampler without it, and the on-manifold discrete sampler, are a compact picture of the whole continuous-sampler predicament. The continuous relaxation buys the ability to run genuine Langevin dynamics, but it pays for that ability by leaving the token manifold, and once off the manifold it faces an impossible choice: either respect the correction and freeze, or ignore the correction and wander through the empty space between tokens where the energy is a poor guide. Neither option samples the intended distribution. The discrete sampler escapes this dilemma by never leaving the manifold, which is why, of the two, it is the one that at least moves sensibly, even though, as the next section shows, the signal it moves by is no better than noise.

\section{The Central Result: The Gradient Fallacy}
\label{sec:results-fallacy}

Everything to this point has concerned the mechanics of the samplers. The result in this section concerns the signal that drives them, and it is the central finding of the thesis. The samplers use the gradient of the frozen likelihood to decide which way to move; the question is whether that gradient points anywhere useful. The ablation of Section~\ref{sec:meth-configs} was designed to answer exactly this, by comparing three proposal methods: the true gradient (\emph{policy}), a random direction rescaled to the gradient's norm (\emph{grad-norm-preserved}), and a fully random vector (\emph{random}). The three-way design is deliberate and is the methodological core of the thesis, because it separates the two things a gradient can contribute. If the gradient's \emph{direction} is informative, the policy method should outperform the norm-preserved random method, which differs from it only in direction. If the gradient's \emph{magnitude} is informative, the norm-preserved method should outperform the fully random method, which differs from it only in magnitude. The expectation, going in, was that the direction at least would matter: that a gradient computed from the model would point toward better tokens more often than a random vector would, even if the magnitude were noisy.

Before the comparison can be trusted, a confound must be removed, and identifying it is itself part of the contribution, because it is the reason a naive version of this experiment would have given an uninterpretable answer. Gradient normalization, which is applied by default for numerical stability, rescales every proposal direction to unit norm. Under normalization, the norm-preserved random method and the fully random method become the same object, a random unit vector, and so an observed equality between the policy method and the random method admits two incompatible explanations: either the direction of the gradient is genuinely uninformative, or the normalization has thrown away a real signal that lived in the magnitude. The normalized ablation cannot tell these apart. The extended ablation therefore repeats the entire comparison \emph{with gradient normalization disabled}, so that the policy and norm-preserved methods both carry the gradient's full unnormalized magnitude and differ from each other only in direction. Any gap between them is then attributable to direction alone, cleanly. This single design decision, doubling the ablation to disentangle direction from magnitude, is what turns a suggestive observation into a claim that survives scrutiny, and it is the reason the study comprises $145$ configurations rather than half as many.

With the confound removed, the result is unambiguous, and it contradicts the expectation. Table~\ref{tab:fallacy} reports the final KL divergence for the policy and norm-preserved methods with normalization disabled and the correction enabled, across four energy functions, together with their difference.
% SOURCE: final_kl_by_model.csv, rows dls.policy.mh.nogn.free.s50 and
% dls.grad_norm_preserved_random_dir.mh.nogn.free.s50, columns gpt2-large / gfn-lb0-500 /
% gfn-lb0-2000 / gfn-lb1-500. policy=(6.474,6.009,6.698,6.008), gnp=(6.430,6.001,6.552,5.338).
The differences are small and, tellingly, inconsistent in sign. On GPT-2 Large the policy method is worse than the random direction by $0.044$; on one GFlowNet variant it is worse by $0.008$; on another it is better by $0.146$; and on the length-normalized variant the random direction is better by a substantial $0.670$. In no case does the true gradient reliably outperform a random direction of the same magnitude. Following the exact gradient is statistically indistinguishable from following a random direction. The direction of the frozen autoregressive gradient carries no usable information for discrete sampling. And because, in the separate comparison, the norm-preserved and fully random methods are themselves indistinguishable, the magnitude carries none either. The gradient, in both its direction and its magnitude, is no better than noise.

\begin{table}[t]
    \centering
    \small
    \begin{tabular}{lccc}
        \toprule
        Energy function & Policy & Grad-norm-preserved & Difference \\
        \midrule
        GPT-2 Large (SFT) & $6.474$ & $6.430$ & $+0.044$ \\
        GFlowNet, \texttt{len\_beta}$=0$, 500 & $6.009$ & $6.001$ & $+0.008$ \\
        GFlowNet, \texttt{len\_beta}$=0$, 2000 & $6.698$ & $6.552$ & $+0.146$ \\
        GFlowNet, \texttt{len\_beta}$=1$, 500 & $6.008$ & $5.338$ & $+0.670$ \\
        \bottomrule
    \end{tabular}
    \caption{Final KL divergence (lower is better) for the true gradient against a norm-preserved random direction, with gradient normalization disabled and the correction enabled, on the discrete sampler at 50 steps. The differences are small and inconsistent in sign, indicating that the direction of the gradient contributes no reliable improvement over a random direction of equal magnitude.}
    \label{tab:fallacy}
\end{table}

There is a single statistically clear contrast in the entire matrix, and it points the wrong way for gradient guidance rather than the right way. On Llama-3, with the correction disabled, the true gradient produces a \emph{higher} final KL divergence than a random direction: there the gradient is not merely uninformative but actively misleading. We defer the mechanistic explanation of this to the next section, but note the shape of it here, because it is instructive: a gradient that points toward the minimizer of a \emph{local linearization} of a globally discontinuous landscape may point somewhere systematically unhelpful, worse than no information at all, if the linearization is a poor guide to the true landscape. That the gradient can be worse than random is a strong hint that the problem is not noise but a systematic mismatch between the local linear model and the true energy, which is exactly what the linearization experiment was built to test. This section has answered the second half of RQ1: the gradient is not a usable search direction, and the ablation locates the failure in the gradient itself rather than in any incidental choice of scale or correction. The obvious next question is why, and it is the subject of the three diagnostic sections that follow.

\subsection{Robustness of the Gradient Fallacy}
\label{sec:results-robust}

Before turning to the mechanism, it is worth recording that the gradient fallacy is not an artifact of a single experimental choice, because a negative result of this centrality invites the suspicion that some incidental setting produced it. Three variations were examined and none overturned it. First, extending the annealing schedule from fifty steps to one hundred does not create a gap between the policy and the random directions: with the longer schedule the policy and grad-norm-preserved methods remain within noise of one another on every model, so the fallacy is not a consequence of insufficient optimization budget.
% SOURCE: final_kl_by_model.csv, s100 rows dls.policy.mh.gn.free.s100 vs
% dls.grad_norm_preserved_random_dir.mh.gn.free.s100: near-identical across all four models.
Second, the phenomenon holds under both settings of the Metropolis--Hastings correction and both settings of gradient normalization, so it is not produced by the interaction of the ablation with any one of those switches. Third, and importantly for the generality of the claim, the same behaviour was observed for multi-token recovery, in which two or three coupled positions are corrupted and must be recovered jointly, as well as for the single-token case reported in the tables. As the number of coupled masked tokens increases the absolute recovery quality degrades, as one would expect from a harder task, but the relative behaviour of the sampling methods is preserved: the policy gradient does not pull ahead of a random direction at any corruption level. The fallacy is therefore a property of the gradient on this class of energy, robust to the schedule, the correction, the normalization, and the difficulty of the recovery task.

\section{Why the Gradient Fails, Part One: The Linearization Radius}
\label{sec:results-linradius}

The gradient fallacy tells us that the gradient carries no information. It does not, by itself, tell us why, and without a mechanism the result is merely a surprising fact rather than an understood one. The linearization experiment, described in Section~\ref{sec:meth-diagnostics}, was built to supply the mechanism by testing the surrogate the discrete proposal actually uses. Recall from Section~\ref{sec:bg-samplers} that the proposal ranks a candidate token $v$ using the first-order term $\hat{\Delta}(v) = \vg^\top(\ve(v) - \ve(x_i))$, a linear approximation to the true change in log-likelihood that a substitution would produce. If the gradient is uninformative, the most direct explanation would be that this surrogate simply does not predict the true change, and the experiment tests that directly by computing both quantities for thousands of candidate substitutions and correlating them.

The result is stark. Figure~\ref{fig:lin-scatter} plots the surrogate against the true change over all candidate--sequence pairs for GPT-2 Large, and there is no relationship whatsoever: the Spearman correlation between the surrogate and the truth is $-0.011$, which is to say zero.
% SOURCE: results_diag/diag_linearization_gpt2sft.json,
% spearman_surrogate_vs_true_ALL=-0.0108, over n_pairs=400000.
Ranking candidate tokens by the gradient surrogate is, quantitatively, equivalent to ranking them at random. This is the gradient fallacy of the previous section, now visible not through the downstream sampling metric but at the level of the proposal itself, which is a much more direct and less confoundable observation. The same measurement on the three GFlowNet variants yields correlations of $0.024$, $0.057$, and $0.046$, all negligible, which confirms that the phenomenon is a property of the frozen-likelihood energy in general and not an accident of the base model.
% SOURCE: diag_linearization_gfn-lb0-2000.json (0.0237), gfn-lb0-500.json (0.0573),
% gfn-lb1-500.json (0.0462).

\begin{figure}[t]
    \centering
    \includegraphics[width=0.82\textwidth]{figures/fig_lin_scatter.pdf}
    % SOURCE: figures/fig_lin_scatter.pdf from plot_diagnostics.py / diag_linearization_gpt2sft.csv
    \caption{The first-order Taylor surrogate used by the discrete proposal, plotted against the true change in sequence log-likelihood, for GPT-2 Large, coloured by the embedding distance of the candidate token. The two quantities are uncorrelated (Spearman $\rho = -0.011$), so ranking tokens by the surrogate is equivalent to ranking them at random.}
    \label{fig:lin-scatter}
\end{figure}

Why should a first-order approximation be this bad? The answer is a mismatch between two length scales, and it is the core mechanism of the thesis. A Taylor expansion is accurate only within a small neighbourhood of the point at which it is taken, and a Transformer \citep{vaswani2017attention}, with its layer normalization and its attention nonlinearities, is a violently nonlinear function of its input embeddings, so the neighbourhood within which its first-order expansion holds is small. The question is whether that neighbourhood is large enough to contain any of the moves the sampler must make. Figure~\ref{fig:lin-radius} answers it by computing the surrogate-to-truth correlation \emph{within} bins of embedding distance, so that one can see how far from the current token the linear model remains valid. The correlation is at best weakly positive at the very smallest distances and has decayed to zero well before the mean inter-token distance of $2.35$ is reached.
% SOURCE: diag_linearization_gpt2sft.json, mean_inter_token_distance=2.354;
% within-bin correlations plotted in figures/fig_lin_radius.pdf.
And the mean inter-token distance is, by its very definition, the scale of the smallest possible substitution: to change a token at all is to move to another token's embedding, and the nearest one is on average this far away. The sampler is therefore always evaluating its linear surrogate at distances far outside the radius within which that surrogate carries any signal. This is the crucial point, and it is worth stating without hedging: the gradient is not merely imprecise here, it is \emph{inapplicable}, because every admissible discrete step leaves the region in which the local linear model means anything. A first-order method is being asked to make finite jumps, and the first order does not survive the jump.

\begin{figure}[t]
    \centering
    \includegraphics[width=0.82\textwidth]{figures/fig_lin_radius.pdf}
    % SOURCE: figures/fig_lin_radius.pdf from plot_diagnostics.py / diag_linearization_gpt2sft.csv
    \caption{Correlation between the gradient surrogate and the true energy change as a function of the embedding distance of the candidate token, for GPT-2 Large. The correlation decays to zero well within the mean inter-token distance (dashed line), so at the distances of any admissible token substitution the surrogate is uninformative. We note that the correlation is low across the whole admissible range rather than high up to a sharp radius and then falling; the term ``linearization radius'' is a summary of this decay, not a claim of a clean threshold.}
    \label{fig:lin-radius}
\end{figure}

A second, independent limitation compounds the first, and the decomposition built into the experiment exposes it. The gradient is taken with respect to the input embedding of the masked position, and it can therefore only register the effect of a substitution on the \emph{future} tokens, whose probabilities depend on that embedding as it flows forward through the network. It is structurally blind to the effect of the substitution on the probability of the candidate token \emph{itself}, given its left context, because that probability depends on the candidate through a discrete index into the output softmax rather than through the continuous input embedding the gradient sees. The decomposition measures both effects. The mean absolute magnitude of the self term is $15.0$ nats and that of the future term is $24.2$ nats, so the self term, to which the gradient is blind, is not a rounding error but a substantial fraction of the total, and the surrogate correlates with the future term ($\rho = 0.03$) rather than the self term ($\rho = -0.10$), exactly as the structural argument predicts.
% SOURCE: diag_linearization_gpt2sft.json: mean_abs_true_delta_self=15.02,
% mean_abs_true_delta_future=24.21, spearman_..._SELF=-0.0956, spearman_..._FUTURE=0.0318.
So the gradient is not only evaluated outside the radius where it would be valid; even within that radius it would be blind to a large part of what determines whether a substitution is good. The full decomposition is shown in the appendix, Figure~\ref{fig:lin-decomp}, and the top-$k$ recall analysis, confirming that gradient ranking does not recover the true best substitutions any better than chance, is shown in Figure~\ref{fig:lin-topk}. Together these explain, at the level of mechanism, the ``why not'' of RQ1, and they explain in particular why the gradient could be worse than random on Llama-3: a surrogate that is not merely uninformative but systematically mis-specified can point actively away from good moves.

\section{Why the Gradient Fails, Part Two: The Likelihood Trap}
\label{sec:results-trap}

The linearization result explains why the samplers cannot move effectively toward low energy. It leaves open a more basic question, one that the failure of the samplers oddly makes it possible to overlook: even if the samplers \emph{could} reach low energy, would low energy be a good place to be? The whole energy-based framework rests on the premise that low energy, that is, high likelihood, corresponds to good text. The likelihood-trap experiment tests that premise directly on the study's own models, and it finds the premise false, reproducing quantitatively on our specific models the neural-text-degeneration phenomenon of \citet{holtzman2020curious}.

The experiment decodes with several strategies and measures the relationship between the per-token likelihood of the generated text and its repetition rate. Figure~\ref{fig:trap} shows the result for GPT-2 Large, and the relationship is monotone and strong: as per-token likelihood increases, that is, as energy decreases, repetition rises and lexical diversity falls. The correlation between per-token likelihood and repetition rate is $+0.65$ on the base model, and it is stronger still on the GFlowNet-tuned variants, reaching $+0.93$ on the length-normalized variant.
% SOURCE: computed from results_diag/diag_likelihood_trap_gpt2sft.csv
% (corr(mean_logp, rep4)=+0.651) and diag_likelihood_trap_gfn-lb1-500.csv (+0.928);
% also gfn-lb0-500 (+0.776), gfn-lb0-2000 (+0.822).
Beam search, which most nearly approximates the true mode of the distribution, produces the highest-likelihood and, at the same time, the most degenerate text of all. To make the degeneration concrete rather than abstract, the following is a representative beam-search output from the length-normalized GFlowNet variant, which achieves among the highest per-token likelihoods in the entire study:

\begin{quote}
\ttfamily aebb-aebb-aebb-aebb-aebb-aebb-aebb-aebb-aebb-aebb-
\end{quote}

This string has a per-token log-probability of $-0.40$, higher than almost any coherent text the model produces, and a four-gram repetition rate of $0.89$.
% SOURCE: diag_likelihood_trap_gfn-lb1-500.csv, beam20 rows: sample
% "aebb-aebb-..." mean_logp=-0.40, rep4=0.89.
It is, by the model's own energy, an excellent sequence, and it is gibberish. This is the likelihood trap in its purest form.

\begin{figure}[t]
    \centering
    \includegraphics[width=\textwidth]{figures/fig_trap_scatter.pdf}
    % SOURCE: figures/fig_trap_scatter.pdf from plot_diagnostics.py / diag_likelihood_trap_gpt2sft.csv
    \caption{The likelihood trap on GPT-2 Large. As mean per-token log-probability increases (energy decreases), the 4-gram repetition rate rises and the distinct-2 diversity falls. The highest-likelihood decoding strategies, greedy and beam search, produce the most degenerate text.}
    \label{fig:trap}
\end{figure}

The implication for the energy-based programme is sharp, and it ties the whole investigation together. Any procedure that succeeds in concentrating probability mass on the lowest-energy regions of this energy, whether an annealed Langevin sampler in its quenching phase or a GFlowNet trained on the likelihood as a reward, is thereby concentrating it on degenerate text. Success at the stated optimization objective is failure at the underlying goal. The gradient-guided samplers, as the earlier sections showed, cannot reach the low-energy regions effectively, and the likelihood trap reveals the strange silver lining in that failure: reaching them would not have helped, because the destination is bad. The methods are, in a sense, protected from their own objective by their inability to optimize it.

The experiment also yields a quantitative fact that the next section needs. The regression of total log-likelihood on generated length, computed over the sampled generations where length varies, has a slope of $-1.12$ nats per token.
% SOURCE: diag_likelihood_trap_gpt2sft.json, length_vs_total_logp.pooled_sampling.slope=-1.123.
Because every additional token lowers the unnormalized log-likelihood by, on average, this much, an energy defined as the raw log-likelihood contains an implicit and substantial incentive toward brevity: shorter sequences are lower-energy simply by virtue of being shorter. This built-in brevity incentive is invisible in the Langevin experiments, which recover fixed-length masked spans, but it becomes decisive the moment a method is free to choose its own length, which is exactly the situation of the GFlowNet. The length-versus-likelihood regression is shown in the appendix, Figure~\ref{fig:trap-length}, and its consequences are the subject of Section~\ref{sec:results-gfn}.

\section{The Geometry Underneath: Embedding Anisotropy}
\label{sec:results-aniso}

Two of the earlier results, the step-size failure of Section~\ref{sec:results-stepsize} and the unreliability of Euclidean distance noted in Section~\ref{sec:meth-metrics}, were promised a fuller geometric explanation, and this section supplies it by reporting the anisotropy analysis in full. The anisotropy of Transformer embeddings is itself a documented phenomenon \citep{gao2019representation, ethayarajh2019contextual}, and measuring it on the two models here both confirms the phenomenon and quantifies the specific consequences it has for the samplers.

In GPT-2 Large the mean pairwise distance between token embeddings is $2.77$ and the mean nearest-neighbour distance is $1.82$; in Llama-3 the corresponding figures are $0.84$ and $0.59$. The mean pairwise cosine similarity is well above zero in both models, at $0.086$ and $0.018$ respectively, which is the signature of the cone-shaped anisotropy: embeddings are not spread through the space but clustered in a narrow region where they point in broadly similar directions.
% SOURCE: diag_anisotropy_gpt2sft.json (mean_pairwise_l2=2.7725,
% mean_nearest_neighbour_l2=1.8222, mean_pairwise_cosine=0.0861) and
% diag_anisotropy_llama3-8b.json (0.8355, 0.5850, 0.0185).
The top principal component of the embedding matrix accounts for only a small fraction of the total variance, so the anisotropy is a broad property of the space rather than the effect of a single dominant direction that could be projected away.
% SOURCE: diag_anisotropy_gpt2sft.json, var_explained_top1=0.0225, var_explained_top10=0.0809.
Figure~\ref{fig:aniso} shows the full distributions of pairwise distances for the two models, on their separate scales.

\begin{figure}[t]
    \centering
    \includegraphics[width=\textwidth]{figures/fig_aniso_hist.pdf}
    % SOURCE: figures/fig_aniso_hist.pdf from plot_diagnostics.py / diag_anisotropy_*_dists.npz
    \caption{Distributions of pairwise token-embedding distances in GPT-2 Large and Llama-3, on separate scales. The mean inter-token distance differs by roughly a factor of three between the two models, which is why no single step-size schedule transfers between them.}
    \label{fig:aniso}
\end{figure}

These measurements close two loops. The roughly threefold difference in inter-token distance between the two models is the direct and quantitative cause of the failure of a Llama-tuned step size on GPT-2 that opened the chapter: a step calibrated to cross a boundary in one geometry is too small to cross one in the other. And the compression of distances that anisotropy produces is the reason a token can be close to the target in embedding space yet grammatically wrong, which is why the study relies on KL divergence, a measure of contextual fit, rather than Euclidean distance as its primary metric. The geometry that trapped the sampler at the start of the investigation is the same geometry that makes the naive metric untrustworthy, and both are consequences of the single documented fact that these embeddings live in a cone.

\section{Consistency Across Models and Architectures}
\label{sec:results-crossmodel}

Before turning to amortization, it is worth pausing to collect the evidence that the phenomena documented so far are consistent across the five energy functions rather than particular to one, because that consistency is what licenses the general claim the discussion will make. The gradient fallacy, measured through the linearization correlation, holds on the GPT-2 base and on all three GFlowNet variants, with correlations of $-0.011$, $0.024$, $0.057$, and $0.046$ respectively, every one of them negligible.
% SOURCE: diag_linearization_*.json spearman_surrogate_vs_true_ALL across the four GPT-2-family models.
The likelihood trap holds on all four, with the likelihood-repetition correlation ranging from $+0.65$ on the base to $+0.93$ on the length-normalized GFlowNet, and it is if anything stronger on the tuned models, which is consistent with tuning having concentrated probability mass more tightly on the degenerate high-likelihood region.
% SOURCE: computed from diag_likelihood_trap_*.csv.
The final-KL values in the gradient-fallacy table move together across models, and the difference between the policy and random directions stays small on every one of them.

The Llama-3 control plays a specific and limited role in this picture, and it is important to be exact about what it does and does not establish. Because Llama-3 uses a different tokenizer and an embedding space roughly a third the scale of GPT-2's, its numbers cannot be placed on the same axis as the GPT-2 family; a KL divergence or an inter-token distance from Llama is not comparable digit-for-digit with one from GPT-2. What the Llama results do establish is qualitative and cross-architectural: the Metropolis--Hastings breakdown occurs on Llama exactly as on GPT-2, with the correction rejecting boundary-crossing moves; the anisotropy is present, though at a different scale; and the gradient is, if anything, more clearly unhelpful, being the one setting in which the true gradient measurably underperforms a random direction. The value of the control is therefore that it rules out the possibility that the entire diagnosis is an artifact of one particular model's idiosyncrasies. The phenomena are not GPT-2 phenomena; they appear in a model of an order of magnitude more parameters, a different tokenizer, and a different embedding geometry, which is what one would expect if their cause is the training objective shared by both rather than any detail of either. This cross-architecture consistency is the last piece of evidence needed before the argument can be made that the cause lies in the energy, which is the subject of the amortization experiment that follows.

\section{Does Amortization Escape the Failure? GFlowNet Fine-Tuning}
\label{sec:results-gfn}

The results to this point form a closed case against the gradient-guided samplers, and they locate the failure in the gradient. This immediately suggests a remedy, and it is the remedy the second half of the thesis was designed to test. If the gradient is the problem, then abandon the gradient. A GFlowNet \citep{bengio2021flow, hu2024amortizing} learns to sample from a reward without ever differentiating it, so it is blind to the linearization failure, blind to the discontinuous drift, blind to everything that defeated the Langevin samplers. The expectation, and it is a reasonable one, is that amortizing the search in this way should escape the pathology. This section reports what happened when the expectation was tested, and the summary is that amortization does not escape the pathology; it relocates it. Three distinct failure modes appeared, and it is worth walking through each, because together they make the diagnosis.

The first is \textbf{capacity collapse}. The initial GFlowNet experiments used the smaller GPT-2 base, and they produced non-linguistic output regardless of how the reward was configured, which indicated that the policy simply lacked the capacity to represent the task at all. Moving to GPT-2 Large, the model of the reference implementation of \citet{hu2024amortizing}, resolved this, and GPT-2 Large is therefore the base for all the GFlowNet results reported here. This mode is worth recording because it is a genuine and separate failure, but it is not the interesting one, because it is a failure of capacity rather than of the energy.

The second is \textbf{length collapse}, and it follows directly and quantitatively from the likelihood trap of the previous section. With the unnormalized reward, at $\texttt{len\_beta}=0$, the reward is the raw log-likelihood, and Section~\ref{sec:results-trap} measured that every additional token lowers that quantity by around $1.12$ nats. The reward therefore contains a strong implicit incentive toward brevity, and the GFlowNet, doing its job faithfully, discovers it: the policy learns to produce extremely short sequences, because short sequences are high-reward simply by being short.
% SOURCE: brevity slope from diag_likelihood_trap_gpt2sft.json (-1.123 nats/token);
% length-collapse behaviour in the free-running generations of the gfn-lb0-* checkpoints.
The extreme form of this is visible in the decoding behaviour: under free-running ancestral sampling, the length-collapse variants never emit an end-of-sequence token at all, running to the generation cap in every case, because the policy has learned that terminating early is the way to keep the reward high and, in the limit, that the safest policy is one that resists stopping in the manner the reward rewards.
% SOURCE: diag_likelihood_trap_gfn-lb0-500.json and gfn-lb0-2000.json,
% frac_emitted_eos=0.0 across all decoding strategies.
The length collapse is not a defect of the optimizer; it is the optimizer correctly following a reward whose implicit incentive the designer did not intend.

The obvious fix is to remove the brevity incentive by normalizing the reward by length, and this produces the third mode, \textbf{reward hacking}. At $\texttt{len\_beta}=1$ the reward is normalized by sequence length, the brevity incentive is gone, and the sequences do indeed return to a normal length. But the output ceases to be linguistic. The policy discovers sequences that achieve high length-normalized likelihood without being coherent text at all, strings of individually high-probability tokens assembled into grammatical nonsense. The following greedy output from this variant is representative:

\begin{quote}
\ttfamily b8d8b8b8b8b8. [\ldots] keredkeredkeredkeredkeredkered
\end{quote}

Such strings recur across the length-normalized variant's high-reward outputs, and they are the reward-hacking counterpart of the length collapse: where the unnormalized reward was gamed by getting shorter, the normalized reward is gamed by finding degenerate high-probability token sequences.
% SOURCE: gfn-lb1-500 generations; diag_likelihood_trap_gfn-lb1-500.csv greedy/beam rows,
% e.g. "b8d8b8b8b8b8...keredkered..." and "aebb-aebb-...".
The two settings of the single length parameter thus produce two different degeneracies, brevity at one extreme and gibberish at the other, and no setting between them was found to produce good text.

These three modes are unified by a single observation, and stating it is the point of the section: the collapse is not fundamentally about length. Length is the visible symptom, the axis along which the failure happens to manifest, but the underlying cause is that the reward, an autoregressive likelihood that was never shaped to serve as an energy over free-form sequences, has no well-behaved optimum where coherent text lies. Every repair of one degeneracy simply surfaces another, because the target itself is malformed. The decisive evidence for locating the cause in the energy rather than in the GFlowNet is a single experiment, and it is the experiment that closes the thesis's argument. If GFlowNet training had reshaped the model's likelihood into a smoother, more navigable landscape, then running the Langevin sampler on the GFlowNet-tuned model should behave differently from running it on the untuned base. It does not. Table~\ref{tab:gfn-unify} reports the final KL divergence of the discrete sampler on the base model and on the three GFlowNet variants, under identical conditions.
% SOURCE: final_kl_by_model.csv, row dls.policy.mh.gn.free.s50:
% gpt2-large=6.541, gfn-lb0-500=6.306, gfn-lb0-2000=6.721, gfn-lb1-500=6.415.
The four values are within noise of one another, and the gradient fallacy replicates on every GFlowNet variant exactly as on the base, as the linearization correlations reported in Section~\ref{sec:results-linradius} already showed. Amortizing the energy does not repair the landscape. The frozen-likelihood pathology survives retraining, which is the strongest evidence the study can offer that the pathology is a property of the energy and not of any sampler or any search procedure. This answers RQ3 directly and in the negative: amortization does not escape the failure, because the failure is in the energy that both the sampler and the amortized policy inherit from the frozen model.

\begin{table}[t]
    \centering
    \small
    \begin{tabular}{lc}
        \toprule
        Energy function & Final KL (DLS, 50 steps, policy, MH, gn) \\
        \midrule
        GPT-2 Large (SFT base) & $6.541$ \\
        GFlowNet, \texttt{len\_beta}$=0$, 500 & $6.306$ \\
        GFlowNet, \texttt{len\_beta}$=0$, 2000 & $6.721$ \\
        GFlowNet, \texttt{len\_beta}$=1$, 500 & $6.415$ \\
        \bottomrule
    \end{tabular}
    \caption{Final KL divergence of the discrete Langevin sampler on the base model and on the three GFlowNet-tuned variants, under identical conditions. The four values lie within noise of one another: Langevin sampling on the amortized energy is indistinguishable from sampling on the untuned base, so the amortization does not repair the energy landscape.}
    \label{tab:gfn-unify}
\end{table}

\section{The Plug-and-Play Claim Tested Directly: Constrained Generation}
\label{sec:results-constrained}

The investigation has established that the fluency energy is unnavigable by gradient and unrepaired by amortization. There remains the question the whole framework was meant to answer, which is RQ4: can an additive constraint steer generation toward a target attribute, as the plug-and-play formulation of \eqref{eq:intro-energy} requires? The constrained-generation extension tests this directly, adding a sentiment constraint to the energy in the manner of \citet{kumar2022gradient} and \citet{qin2022cold} and measuring whether the target sentiment can be induced. Given everything that precedes it, the expectation is pessimistic, but the test is worth running directly rather than inferring its outcome, because the constraint gradient is a different object from the fluency gradient and might conceivably behave better.

It does not behave better in a way that produces control. Table~\ref{tab:constrained} reports the steering effect for the discrete sampler across the five constraint modes on the infilling task with a positive-sentiment target.
% SOURCE: results_constrained/gpt2-large.infill.ours.dls.init-random_token.*.lbl1.json,
% steering_gain: lm_only=0.0, cons_only=-8.0, full=-1.0, random=-7.0, cons_random=+1.0;
% final_kl: 8.558, 8.308, 8.541, 8.512, 8.528.
The steering gains are small in magnitude and inconsistent in sign. The full combined energy, which is the intended plug-and-play configuration, produces a gain of $-1.0$, essentially no steering and if anything in the wrong direction; the constraint-only mode produces $-8.0$, worse than the baseline; and no mode produces the systematic movement toward the target sentiment that a working constraint would require. This is the outcome the preceding sections predict: the constraint gradient is subject to the same discrete-manifold difficulties as the fluency gradient, and adding it to an energy landscape that is already unnavigable does not produce controllable generation. Inspection of the generated text confirms the picture, showing the same degeneracy observed throughout; a representative infilled span reads \texttt{The rostr Bill of H PM americanus bears one or connected sp}, in which the masked positions have been filled with locally plausible but globally incoherent tokens.
% SOURCE: results_constrained/gpt2-large.infill.ours.dls.init-random_token.full.lbl1.json,
% examples field, recovered span text.
The plug-and-play promise of \eqref{eq:intro-energy}, that an additive constraint can steer a frozen model at inference time, is not realized on this energy landscape, because its precondition, a navigable fluency energy, does not hold. This answers RQ4 in the negative, and, importantly, it does so for a reason the earlier sections have already established rather than as an isolated and unexplained failure, which is what makes it evidence for the thesis's single underlying diagnosis rather than merely one more disappointing table.

\begin{table}[t]
    \centering
    \small
    \begin{tabular}{lcc}
        \toprule
        Constraint mode & Steering gain & Final KL \\
        \midrule
        lm\_only & $0.0$ & $8.558$ \\
        cons\_only & $-8.0$ & $8.308$ \\
        full & $-1.0$ & $8.541$ \\
        random & $-7.0$ & $8.512$ \\
        cons\_random & $+1.0$ & $8.528$ \\
        \bottomrule
    \end{tabular}
    \caption{Steering effect of the sentiment constraint under the discrete sampler on the infilling task, positive-sentiment target. The steering gain is the change in the percentage of generations classified as the target sentiment relative to the unconstrained baseline. The effects are small and inconsistent in sign; no mode produces reliable steering toward the target.}
    \label{tab:constrained}
\end{table}
